\chapter{Middleware and Observability}
\label{ch:observability}

\begin{quote}\itshape
Agent Framework will tell you what happened. It will not tell you whether it was
any good.
\end{quote}

\noindent
The previous three chapters built systems with a lot of moving parts: graphs that
drop messages silently, runs that resume in a different process, five agents
handing work between themselves. Every one of those chapters ended by pointing
here, because none of them can be operated without seeing inside them.

This chapter covers the two mechanisms that make that possible --- middleware,
which is where you insert your own behaviour, and OpenTelemetry, which is where
the framework tells you what it did. Both are good. Then it covers the gap
between them, which is the thesis of everything that follows: the framework
instruments cost and latency thoroughly and quality not at all, and quality is
the thing you actually needed to know.

\begin{verifybox}
Derived from the source of \pkg{Microsoft.Agents.AI} --- the agent builder,
the OpenTelemetry decorator and the compaction telemetry --- and from the
\code{Observability/} folder of \pkg{Microsoft.Agents.AI.Workflows}. Activity
source names, span names and tag keys are transcribed from the constants in that
source and are reliable. The scoring middleware in \S\ref{sec:obs-scoring} is the
author's own code and has not been compiled.
\end{verifybox}

% ============================================================
\section{The middleware pipeline}
\label{sec:obs-middleware}
\index{AIAgentBuilder@\texttt{AIAgentBuilder}}
\index{DelegatingAIAgent@\texttt{DelegatingAIAgent}}
\index{middleware!agent level}
\index{middleware!function level}
\index{middleware!chat client level}
\index{FunctionInvokingChatClient@\texttt{FunctionInvokingChatClient}}
\index{middleware}

There is no separate middleware abstraction to learn. Middleware in this
framework is the decorator pattern with a builder in front of it: an agent wraps
an agent, and the builder composes the wrapping. If you have written ASP.NET Core
middleware the shape will be familiar, and if you have used
\code{IChatClient} pipelines it is the same idea one level up.

\begin{csharp}[caption={The builder}, label={lst:obs-builder}]
AIAgent agent = new AIAgentBuilder(innerAgent)
    .UseOpenTelemetry()
    .UseLogging(loggerFactory)
    .Build();
\end{csharp}

Three places accept your code, and choosing the wrong one is the most common
mistake in this area --- not because it fails loudly, but because the thing you
wanted to observe happens at a different level from the one you hooked.

\begin{center}
\small
\begin{tabularx}{\linewidth}{@{}lX@{}}
\toprule
\textbf{Level} & \textbf{Sees} \\
\midrule
Agent & One agent run: the messages in, the response out \\[4pt]
Function & Each tool call the agent makes, including arguments and result \\[4pt]
Chat client & Every model request, including each iteration of the tool loop \\
\bottomrule
\end{tabularx}
\end{center}

\subsection{Agent-level}

The outermost level. One invocation in, one response out --- regardless of how
many model calls and tool calls happened inside.

\begin{csharp}[caption={Agent middleware from a delegate}, label={lst:obs-agent-mw}]
AIAgent agent = new AIAgentBuilder(inner)
    .Use(async (messages, session, options, next, cancellationToken) =>
    {
        Stopwatch sw = Stopwatch.StartNew();
        await next(messages, session, options, cancellationToken);
        this._metrics.RecordTurn(sw.Elapsed);
    })
    .Build();
\end{csharp}

There are richer overloads --- one that hands you the inner agent and lets you
return a response without calling it at all, which is how caching and short
circuiting are written, and one taking separate delegates for the streaming and
non-streaming paths. For anything beyond a few lines, derive from
\api{DelegatingAIAgent} instead and override what you need; it is the same
mechanism with somewhere to put fields.

\begin{warning}
\textbf{Implement both paths or neither.} An agent has a streaming and a
non-streaming entry point, and middleware supplied for only one leaves the other
uninstrumented. A quality gate that runs on \code{RunAsync} and silently does not
run on \code{RunStreamingAsync} is worse than no gate, because it is believed.
\end{warning}

\subsection{Function-level}

One level in: each tool invocation, with its arguments and its result. This is
where approval gates, argument validation and per-tool timing belong.

\begin{csharp}[caption={Intercepting tool calls}, label={lst:obs-function-mw}]
AIAgent agent = new AIAgentBuilder(inner)
    .Use(async (agent, context, next, cancellationToken) =>
    {
        if (IsDestructive(context.Function.Name) && !await ApprovedAsync(context))
        {
            return "Refused: this tool requires approval.";
        }

        return await next(context, cancellationToken);
    })
    .Build();
\end{csharp}

\begin{warning}
\textbf{Function middleware requires a function-invoking chat client somewhere in
the pipeline.} If the inner agent's pipeline does not include one, the agent this
adds will throw when it is invoked --- not when it is built. The failure is at run
time, on the first tool call, in whatever environment first exercised a tool.

If you are constructing the chat client by hand rather than through the usual
agent construction path, this is the thing to check first.
\end{warning}

\subsection{Chat-client-level}

The innermost level, and the one people reach for last and should often reach for
first. An agent run that makes six model calls in its tool loop is one agent-level
invocation and six chat-client-level ones.

That difference is exactly the compaction point from Chapter~\ref{ch:memory}:
anything that needs to act on \emph{every} model request --- token accounting,
context compaction, request-level caching --- belongs here, because at the agent
level the tool loop has already happened and the tokens are already spent.

\subsection{Exception handling in the pipeline}

Middleware wraps a call, so a \code{try}/\code{catch} around \code{next} does what
you expect. Three things are worth being deliberate about.

\textbf{Decide whether you are observing or intervening.} Middleware that catches,
records and rethrows is observation and is safe. Middleware that catches and
returns a substitute response has changed the agent's behaviour, and the caller
cannot tell. Both are legitimate; conflating them is not.

\textbf{Streaming exceptions arrive late.} On the streaming path the exception
surfaces while the caller enumerates, not when the method is called. A
\code{try}/\code{catch} around the call to \code{next} that returns an enumerable
will not catch it.

\textbf{Never swallow cancellation.} A cancelled run that is reported as a
completed one produces a trace that says everything was fine, which is the worst
possible output from an observability layer.

% ============================================================
\section{OpenTelemetry out of the box}
\label{sec:obs-otel}
\index{activity source}
\index{OpenTelemetry!GenAI conventions}
\index{observability!sensitive data}
\index{WorkflowTelemetryOptions@\texttt{WorkflowTelemetryOptions}}
\index{span!workflow}
\index{OpenTelemetry}

Tracing is a decorator you add to the pipeline. It implements the OpenTelemetry
semantic conventions for generative AI systems, so the spans are the standard
ones rather than a Microsoft dialect, and any conformant backend will understand
them.

\begin{csharp}[caption={Turning on tracing}, label={lst:obs-otel-agent}]
AIAgent agent = new AIAgentBuilder(inner)
    .UseOpenTelemetry()
    .Build();
\end{csharp}

Workflows are instrumented separately, on the builder:

\begin{csharp}[caption={Tracing a workflow}, label={lst:obs-otel-workflow}]
Workflow workflow = new WorkflowBuilder(ingest)
    .AddEdge(ingest, classify)
    .WithOutputFrom(classify)
    .WithOpenTelemetry()
    .Build();
\end{csharp}

\subsection{Two activity sources, and you must register both}

This is the detail that costs people an afternoon. Agents and workflows emit
under \emph{different} activity source names:

\begin{center}
\small
\begin{tabularx}{\linewidth}{@{}lX@{}}
\toprule
\textbf{Emitted by} & \textbf{Source name} \\
\midrule
Agents, chat clients, compaction & \code{Experimental.Microsoft.Agents.AI} \\
Workflows & \code{Microsoft.Agents.AI.Workflows} \\
\bottomrule
\end{tabularx}
\end{center}

\noindent
Register one and you get half a trace, with no error to tell you the other half
exists. The symptom is a workflow trace whose executor spans contain no model
calls, or a set of agent spans with no graph around them.

\begin{versionbox}
Note the \code{Experimental.} prefix on the agent source name. It is not
decoration --- it is the framework telling you that the source name is not yet
stable. Put it in configuration rather than in a string literal scattered through
your startup code, because it will change, and when it does your dashboards will
go quiet rather than break.
\end{versionbox}

\subsection{Which spans you get}

Workflow instrumentation produces these, and the tag keys are worth knowing
because \S\ref{sec:obs-reading} debugs with them:

\begin{center}
\small
\begin{tabularx}{\linewidth}{@{}lX@{}}
\toprule
\textbf{Span} & \textbf{Covers} \\
\midrule
\code{workflow.build} & Graph construction and validation \\
\code{workflow.session} & A run \\
\code{workflow\_invoke} & An invocation of the workflow \\
\code{executor.process} & One executor handling one message \\
\code{edge\_group.process} & One delivery attempt along an edge \\
\code{message.send} & A message emission \\
\bottomrule
\end{tabularx}
\end{center}

\noindent
Agent spans follow the generative-AI conventions, with an operation named for
agent invocation and the usual attributes --- agent id, agent name, agent
description, provider name --- alongside the token counts the conventions define.

\subsection{Compaction telemetry}
\index{compaction!telemetry}

Chapter~\ref{ch:memory} promised this. Compaction emits its own activities ---
one for the compaction itself, one for the provider invocation, one for the
summarisation path --- and the tags are the measurement that section asked for:

\begin{center}
\small
\begin{tabularx}{\linewidth}{@{}lX@{}}
\toprule
\textbf{Tag} & \textbf{Records} \\
\midrule
\code{compaction.strategy} & Which strategy ran \\
\code{compaction.triggered} & Whether the threshold was crossed \\
\code{compaction.compacted} & Whether anything was actually removed \\
\code{compaction.before.tokens} & Token count before \\
\code{compaction.after.tokens} & Token count after \\
\code{compaction.before.messages} & Message count before \\
\code{compaction.after.messages} & Message count after \\
\code{compaction.duration\_ms} & What it cost in time \\
\code{compaction.groups\_summarized} & Groups summarised, on that path \\
\code{compaction.summary\_length} & Length of the summary produced \\
\bottomrule
\end{tabularx}
\end{center}

\begin{note}
The pair to watch is \code{triggered} against \code{compacted}. A strategy that
triggers frequently and compacts nothing is doing work and buying nothing, and it
is invisible in any dashboard that only charts token counts. That comparison is
the cheapest useful alert in this whole chapter.
\end{note}

\subsection{Sensitive data}

Both the agent decorator and the workflow options have a switch that puts message
content --- prompts, responses, executor inputs and outputs --- into span
attributes. It is off by default.

\begin{warning}
\index{observability!PII}
Turning it on sends user input to your tracing backend, where it is retained
under that backend's policy rather than your application's, is visible to
everyone with dashboard access, and is almost certainly outside whatever your
privacy notice describes. That is a defensible trade in local development and
rarely one in production.

If you need it in production, need it \emph{narrowly}: a sampled fraction, a
redaction step in front of it, and a retention policy you have actually
configured. ``We turned on sensitive data to debug an incident and left it on''
is a common story with a bad ending.
\end{warning}

The workflow options also let you disable individual span kinds --- build, run,
executor processing, edge processing, message sends. Edge and message spans are
by far the highest volume, and a graph running at any real rate will produce more
of them than anything else in your system. Turning those two off while keeping
executor spans is a reasonable production posture, with the significant caveat
that it disables exactly the spans \S\ref{sec:obs-reading} needs.

% ============================================================
\section{Where traces go}
\label{sec:obs-destinations}
\index{Application Insights}
\index{OTLP}
\index{sampling}

Nothing here is agent-specific, which is the point --- these are ordinary .NET
activities and every standard destination works.

Locally, the Aspire dashboard is the fastest thing to stand up and the best trace
viewer of the three; it runs in a container, speaks OTLP, and needs no account.
For a hosted application, Application Insights is the path of least resistance on
Azure. Anything OTLP-compatible works: an OpenTelemetry Collector in front of
whatever you already run is the arrangement that avoids a decision.

\begin{csharp}[caption={Wiring up tracing}, label={lst:obs-wireup}]
builder.Services.AddOpenTelemetry()
    .WithTracing(tracing => tracing
        .AddSource("Experimental.Microsoft.Agents.AI")   // agents
        .AddSource("Microsoft.Agents.AI.Workflows")      // workflows
        .AddOtlpExporter());
\end{csharp}

\begin{note}
Sampling deserves a decision rather than a default. Agent runs are low-volume and
expensive compared with HTTP requests, so the usual instinct to sample at one per
cent is wrong here --- you will lose the rare failing run, which is the only one
worth having. Trace every agent run and sample the edge and message spans
underneath them if volume becomes a problem.
\end{note}

% ============================================================
\section{Reading a multi-agent trace}
\label{sec:obs-reading}
\index{span!links}
\index{span!siblings not nested}
\index{observability!trace shape}
\index{observability!reading traces}

\subsection{The shape is not what you expect}

A trace of a single agent is a waterfall: an agent span, model calls nested
inside it, tool calls nested inside those. Familiar, and easy to read.

A workflow trace is not that, and this catches everyone. \textbf{Executor
processing spans are siblings, not children of one another.} The framework's own
source says so explicitly, and explains why: the causal relationship between one
executor and the next is represented with span \emph{links} rather than with
parent-child nesting.

Once stated it is obviously right. Chapter~\ref{ch:workflows-basics} established
that a superstep runs its executors independently and delivers their messages to
the next superstep --- so executor A does not \emph{contain} executor B, it
precedes it. Nesting would misrepresent a fan-out, where three executors genuinely
run at once, as a hierarchy.

\begin{warning}
Most trace viewers render parent-child nesting beautifully and span links poorly
--- often as a small icon you have to know to click. So the default view of a
workflow trace is a flat list of executor spans with the causality hidden.

The consequence is practical: \textbf{do not try to read workflow causality off
the waterfall}. Read it off the edge spans, which is what the next subsection
does.
\end{warning}

\needscreenshot{aspire-dashboard-trace-waterfall}{A multi-agent run, as the trace
shows it}{The Aspire dashboard trace view for a workflow run containing a fan-out
to two agents and a fan-in. Capture the full span list with the workflow session
span, the sibling executor spans, the edge spans between them, and the nested
model-call spans under the agent executors. The point is the contrast between
nesting under an agent and sibling ordering between executors, so both must be
visible in one frame. Aspire dashboard only, not Visual Studio.}

\subsection{Finding the dropped delivery}
\label{sec:obs-dropped}
\index{edge!delivery status tag}
\index{observability!query}

Chapter~\ref{ch:workflows-basics} \S\ref{sec:wf-mismatch} left an IOU: a workflow
that produces no output is usually a silently dropped message, and the way to
find it is in the trace. Here is the procedure.

Every delivery attempt produces an \code{edge\_group.process} span, and that span
carries two tags: a boolean recording whether the message was delivered, and a
string recording why not if it was not. The values of the second are those from
\S\ref{sec:wf-mismatch} --- \code{delivered}, \code{dropped type mismatch},
\code{dropped target mismatch}, \code{dropped condition false}, \code{buffered},
\code{exception}.

So the query is a filter, and it is the single most valuable thing in this
chapter:

\begin{shellcmd}[caption={The query to save}, label={lst:obs-query}]
# Every delivery in this run that did not happen, and why.
span.name = "edge_group.process" AND edge_group.delivered = false
\end{shellcmd}

Run that against a trace of a workflow that produced nothing and you will
normally have the answer in one step. The span's source and target tags name the
edge; the status tag names the cause; and the cause tells you which fix applies:

\begin{center}
\small
\begin{tabularx}{\linewidth}{@{}lX@{}}
\toprule
\textbf{Status} & \textbf{What to change} \\
\midrule
\code{dropped type mismatch} & The target's declared input type. Usually an
interface where a concrete type was needed \\[4pt]
\code{dropped condition false} & The edge predicate, or the value it tests \\[4pt]
\code{dropped target mismatch} & The addressed target id \\[4pt]
\code{buffered} & Nothing --- a barrier is still waiting. If it never releases,
one of its sources is dropping \\
\bottomrule
\end{tabularx}
\end{center}

\begin{note}
Save that filter as a dashboard the day you deploy your first workflow, before
you need it. The failure it catches is silent, which means nobody goes looking
for it --- they go looking for why the output is empty, which is a much longer
path to the same span.
\end{note}

% ============================================================
\section{DevUI}
\label{sec:obs-devui}
\index{DevUI!loopback default}
\index{DevUI!authentication}
\index{DevUI}

Traces tell you what happened after the fact. \pkg{Microsoft.Agents.AI.DevUI} is
for the other half of the job: poking at an agent while it runs. It is a web UI
you mount into your own application, so it inspects the real agent with the real
configuration rather than a copy.

\begin{csharp}[caption={Mounting DevUI}, label={lst:obs-devui}]
builder.AddDevUI();

WebApplication app = builder.Build();

app.MapDevUI();
\end{csharp}

\begin{note}
\textbf{DevUI is loopback-only by default,} and it warns at startup if you
configure it otherwise. Remote access is an explicit opt-in, and there is bearer
token authentication --- configurable in options or through an environment
variable --- for when you need it reachable from elsewhere.

That default is the right one. A UI that can invoke your agents with your
credentials is not something to expose casually, and the fact that it is
convenient to reach from a colleague's machine is precisely the reason it ends up
reachable from more than that.
\end{note}

\needscreenshot{devui-agent-inspect}{Inspecting a live agent}{The DevUI interface
with an agent selected, showing a conversation in progress and the inspection
pane with the agent's tools and configuration visible. Include at least one tool
call expanded so its arguments and result can be read. Browser only, no IDE.}

There is also an Aspire integration, which is worth knowing about if you are
already using Aspire: it wires the developer UI into the app host alongside the
dashboard, so traces and interactive inspection sit behind one launch.

% ============================================================
\section{What the traces do not measure}
\label{sec:obs-gap}
\index{quality!not instrumented}
\index{observability!what it misses}

Everything so far works. The instrumentation is thorough, standards-based and
close to free. Now the honest part.

Read back what is actually recorded. Token counts. Durations. Which executor ran,
which edge delivered, which model was called, whether an exception was thrown.
Every one of those is a fact about \emph{mechanism}. Not one of them is a fact
about whether the agent did a good job.

\begin{center}
\small
\begin{tabularx}{\linewidth}{@{}lX@{}}
\toprule
\textbf{Instrumented} & \textbf{Not instrumented} \\
\midrule
Tokens consumed & Whether the answer was correct \\
Latency & Whether it answered the question asked \\
Which tools were called & Whether calling them was the right choice \\
Turn count & Whether the turns made progress \\
Errors thrown & Whether a successful run was useless \\
\bottomrule
\end{tabularx}
\end{center}

This is not a criticism of the framework. Correctness is domain-specific and no
framework can know it. But it has a consequence that is easy to miss and
expensive to learn: \textbf{a green dashboard is compatible with an agent that is
failing every user}. Latency flat, error rate zero, token cost within budget, and
every answer wrong. Nothing in the box will page you.

\begin{note}
Magentic's progress ledger from Chapter~\ref{ch:orchestration} is the one place
the framework goes further --- it explicitly asks whether progress is being made
rather than whether work is being done. It is worth noticing that this exists in
one orchestration pattern and nowhere else, and that it needed a language model
to compute it. That is a clue about what closing this gap costs.
\end{note}

This gap is the reason Chapter~\ref{ch:evaluation} exists, and the reason
Chapter~\ref{ch:capstone} builds what it builds. The rest of this chapter takes
the first step: if quality is not instrumented, instrument it yourself.

% ============================================================
\section{Custom middleware that scores a turn}
\label{sec:obs-scoring}
\index{scoring!as span attribute}
\index{middleware!scoring}
\index{refusal detection}

The mechanism is already in place. Middleware sees the messages in and the
response out; the ambient activity is the span the framework created. So a score
computed in middleware and set as a tag lands on the agent's own span, next to
the token counts, queryable with the same tools.

\begin{verifybox}
The author's own code, uncompiled. The shape is verified --- middleware
registration, and setting tags on the current activity, are both standard --- but
treat it as a design to type in.
\end{verifybox}

\begin{csharp}[caption={A quality score as a span attribute}, label={lst:obs-score}]
AIAgent agent = new AIAgentBuilder(inner)
    .UseOpenTelemetry()   // outermost: the span must exist before we tag it
    .Use(async (messages, session, options, innerAgent, cancellationToken) =>
    {
        AgentResponse response =
            await innerAgent.RunAsync(messages, session, options, cancellationToken);

        TurnScore score = await this._scorer.ScoreAsync(messages, response, cancellationToken);

        Activity? activity = Activity.Current;
        activity?.SetTag("quality.score", score.Value);
        activity?.SetTag("quality.scorer", score.ScorerName);
        activity?.SetTag("quality.refused", score.WasRefusal);

        return response;
    }, runStreamingFunc: null)
    .Build();
\end{csharp}

\begin{warning}
\textbf{Order matters and is easy to get backwards.} The tracing decorator must be
\emph{outside} the scoring middleware, so that the span is open when the score is
computed. Reverse them and \code{Activity.Current} is whatever ambient span the
caller had --- often the incoming HTTP request --- and your scores will attach to
the wrong thing without ever being null.

Assert on this in a test: run once and confirm the tag is on a span named for
agent invocation, not on the request span.
\end{warning}

\subsection{What to score, given that you cannot score correctness}

The scorer is the hard part, and it is worth being clear about what is cheap and
what is not.

\textbf{Cheap and immediately useful:} refusal detection, empty or truncated
responses, whether the response cites a tool result it was given, response length
outside expected bounds, and whether a tool loop terminated or hit its cap. None
of these needs a model, all of them catch real failures, and together they cover
a surprising fraction of ``the agent stopped working'' incidents.

\textbf{Expensive:} a model-graded score. It is another model call per turn, on
the hot path, with its own latency and failure modes --- and its own quality
problem, one level up. If you do it, do it out of band on a sample, not inline on
every turn.

\begin{note}
Even a crude score changes what you can ask. ``Did the refusal rate move after
yesterday's prompt change'' is answerable with a boolean tag and a dashboard, and
unanswerable without one. That question --- \emph{did the change make things
worse} --- is Chapter~\ref{ch:evaluation}'s subject, and a span tag is the
cheapest possible first version of the answer.
\end{note}

\begin{exercisebox}
Before Chapter~\ref{ch:evaluation}:

\begin{enumerate}[leftmargin=1.4em]
\item Register only the agent activity source, run a workflow, and look at the
      trace. Then add the workflow source. The difference is what
      \S\ref{sec:obs-otel} is warning you about, and seeing it once is worth more
      than reading it.
\item Break an edge's types on purpose, as Chapter~\ref{ch:workflows-basics}
      asked, and this time find it with the query in \S\ref{sec:obs-dropped}.
      Time yourself. Then find the same bug without the trace.
\item Instrument the same tool-heavy run with middleware at all three levels and
      count how many times each fires. The ratio is the tool loop, and it is
      usually larger than people expect.
\item Add the refusal-detection scorer. Deploy a deliberately over-cautious
      system prompt and confirm the dashboard shows it. This is a quality
      regression caught by an operational signal, which is the whole idea.
\item Turn on sensitive data in a local Aspire run, look at what is now in your
      spans, and decide honestly whether you would ship that configuration.
\end{enumerate}
\end{exercisebox}

% ============================================================
\section{Summary}

Middleware is the decorator pattern with a builder in front of it, and it comes
at three levels. Agent-level sees one run; function-level sees each tool call;
chat-client-level sees every model request including each iteration of the tool
loop. Anything that must act on every model request belongs at the innermost
level, because by the time the agent level sees the run, the tokens are spent.
Implement the streaming and non-streaming paths together, or leave one silently
uninstrumented.

Tracing is a decorator on agents and an extension on the workflow builder, and it
emits the standard generative-AI semantic conventions. Agents and workflows use
two different activity source names; register both or you get half a trace with
nothing to tell you the other half exists. The agent source is still prefixed as
experimental, which is a promise that it will change.

Workflow executor spans are siblings joined by links, not nested children,
because a superstep runs its executors independently. Most viewers render links
poorly, so do not read causality off the waterfall --- read it off the edge spans.
Every delivery attempt is tagged with whether it was delivered and why not, which
turns Chapter~\ref{ch:workflows-basics}'s silent drop from an afternoon of
bisection into one filter.

Compaction emits its own before-and-after counts, and the pair worth alerting on
is triggered against compacted: a strategy that fires often and removes nothing
costs money and shows up nowhere else.

Everything instrumented is a fact about mechanism. Tokens, latency, which
executor ran, which tool was called, whether an exception was thrown. Nothing in
the box records whether the answer was any good, which means a green dashboard is
compatible with an agent failing every user. Middleware plus a span tag is the
cheapest first step towards closing that --- refusal detection and truncation
checks cost nothing and catch real incidents.

Chapter~\ref{ch:evaluation} does the job properly.
